\section{Introduction}
\label{sec:intro}

Pedestrian trajectory forecasting takes the recent positions of the people in a scene and predicts where they will move over the next few seconds, typically as a set of candidate futures with associated probabilities. It is a core component of autonomous driving, mobile robot navigation and crowd management, where the forecast is consumed by a downstream planner that must reason about the intentions of other agents. A decade of work has produced increasingly expressive models of social interaction and of the multi-modality of human motion, progressing from social pooling over neighbouring trajectories \cite{alahi2016social,gupta2018social}, through graph and attention based scene encoders \cite{mohamed2020social,yuan2021agentformer,liang2020lane,gao2020vectornet}, to generative decoders that represent the future as a distribution \cite{mangalam2020not,salzmann2020trajectron,fu2025moflow,zhang2024demo}.

Almost all of this work shares an assumption that is rarely stated: the history is complete, contiguous and regularly sampled. The input is a tensor of shape $T\times 2$ in which every row is a measurement. Real systems do not provide such tensors. Histories are assembled by a detector and a multi-object tracker, and they inherit every failure of that pipeline: a pedestrian walks behind a parked car and the detector loses them for four frames; two people cross and the tracker swaps their identities; the camera drops frames under load. The result is a history in which the \emph{set} of observed frames differs across agents and across episodes, and in which a missing row carries no positional information at all.

The standard response is to make the history complete before forecasting. A statistical or learned imputer fills the gaps, and an off-the-shelf predictor consumes the repaired sequence \cite{chib2024pedestrian,xu2023uncovering}. A more recent line couples the two stages and optimises them jointly \cite{chib2024mstip,chib2026hypergraph,sun2025rmpyolo}. Both designs inherit a structural problem. Imputation is trained to minimise reconstruction error on the observed data distribution, whereas forecasting is scored by displacement error on the future; the two objectives coincide only when the reconstructed value is the conditional mean of the true position, which is precisely not the case when the gap is long or the motion is multimodal. Worse, once the gap has been filled, the predictor can no longer tell which rows were measured and which were inferred, so a hallucinated coordinate propagates into velocity estimates, into the social relations computed against neighbours, and finally into the forecast, with the same authority as a measurement.

An alternative is to let the predictor consume the incomplete history directly and to model the \emph{evidence} available at each frame as an explicit quantity. This reframing has three consequences that the imputation view obscures. First, the reliability of a history state is not a property of a single frame but of a sequence: it decays with the age of the last observation, and it should modulate how strongly that state enters a temporal recurrence. Second, missingness changes the geometry of social interaction: an agent that is currently unobserved may still query its neighbours, but its stale position must not be broadcast to them as if it were current. Third, the decoder's starting point is itself uncertain: when the last frames are missing, the most recent trustworthy position is older than the forecast origin, and the decoder must be told how old it is rather than being anchored at a placeholder row.

Realising this programme requires two things that are currently missing from the literature. The first is a \emph{protocol}. Published studies on missing histories disagree on how missingness is generated, how samples are split, which checkpoint is reported and which metric defines the ranking; some report imputation-then-prediction and direct prediction in the same table even though the two settings expose the model to different information. Numbers are therefore not comparable across papers, and the effect of a single design choice cannot be isolated. The second is a \emph{mechanism} that carries the evidence state through the whole model. Existing mask-aware methods typically concatenate a binary mask to the input features \cite{chen2023trajmae,yang2023rmp,chib2023dropping}, or use masking as a pre-training objective, or restrict themselves to shortening the history uniformly \cite{sun2022momentary,monti2022howmany,qiu2025adapting}. None of them makes the evidence state a first-class variable that simultaneously governs which agents may act as spatial evidence, how strongly each frame enters the temporal recursion, and where the decoder starts.

This paper addresses both gaps for the third-person (fixed-camera or bird's-eye) setting. We first freeze \textbf{TrajGap-Bench}, a benchmark built on the public TrajImpute release \cite{chib2024pedestrian}. Rather than publishing new trajectories, the benchmark fixes the experimental contract: the Easy and Hard variants of the release are merged at sample level into a single Mixed split for training, validation and test; a single checkpoint rule based on the validation $\mathrm{minFDE}_{20}$ of the final prediction head is applied to every method; only two metrics, $\mathrm{minADE}_{20}$ and $\mathrm{minFDE}_{20}$ with $K=20$, define the ranking; and five diagnostic dimensions (missing-frame count, anchor lag, longest missing run, number of valid frames, terminal missingness) are reported as stratification rather than as additional leaderboards. We then propose \textbf{GAST} (Gap-Aware Spatio-Temporal forecaster), which represents the history of every agent by an observation-evidence state, uses it to gate a frame-wise relational encoder, modulates a causal state-space temporal encoder and its readout, and conditions a goal-based multimodal decoder on an anchor state that records how far the last observation is from the forecast origin. The model never imputes a coordinate, and when no frame is missing its evidence conditioning becomes an identity, so it falls back to a complete-history backbone.

We evaluate GAST against six external baselines on all five scenes of the benchmark, comprising 302{,}886 test samples. The main findings are as follows.
\begin{itemize}
\item GAST attains the best macro-averaged $\mathrm{minADE}_{20}$/$\mathrm{minFDE}_{20}$ among all compared methods (0.3085/0.4716~m), 11.8\%/9.8\% below the strongest external baseline (MoFlow). It is better than MoFlow on all five scenes for $\mathrm{minADE}_{20}$ and on four for $\mathrm{minFDE}_{20}$.
\item The advantage is severity-dependent rather than uniform. Relative to MoFlow, GAST's margin grows monotonically with the number of missing frames, from $-7.1\%$ at zero missing to $-22.3\%$ at six, before narrowing at the degenerate setting in which only one frame survives. This pattern is the main evidence that the gain comes from evidence routing rather than from extra capacity.
\item Two negative results are reported in full. The improvement is not uniform across scenes (GAST trails AgentFormer on HOTEL-M and MoFlow slightly on ZARA1-M), and the best-of-20 gains are accompanied by a small top-1 ranking deficit relative to MoFlow ($\mathrm{FDE@1}$ $+4.1\%$). We analyse both and argue that they constrain how the method should be used and extended.
\end{itemize}

The remainder of the paper is organised as follows. Section~\ref{sec:related} reviews related work. Section~\ref{sec:bench} defines the task and the benchmark protocol. Section~\ref{sec:method} presents GAST. Section~\ref{sec:exp} reports experiments and ablations, Section~\ref{sec:disc} discusses the findings and their limitations, and Section~\ref{sec:concl} concludes.
